\section{The Chebyshev filter as a spectral projector}

Let $H$ be a self-adjoint operator on a finite-dimensional inner product space,
with orthonormal eigenbasis $b_j$ and eigenvalues $h_j$.

\begin{theorem}[Polynomials of an operator on its eigenvectors]
    \label{thm:polynomial_on_eigenvector}
    \lean{ContinuousLinearMap.aeval_map_apply_of_apply_eq_smul}
    \leanok
    If $Hv=\mu v$ with $\mu$ real and $p$ is a real polynomial, then
    $p(H)v=p(\mu)v$.
\end{theorem}

\begin{proof}
    \leanok
    Both sides are additive in $p$, and $H^mv=\mu^mv$ by induction on $m$.
\end{proof}

\begin{theorem}[A polynomial filter as a spectral projector]
    \label{thm:polynomial_spectral_projector}
    \lean{Submodule.isSpectralProjector_aeval}
    \leanok
    \uses{thm:polynomial_on_eigenvector, def:spectral_projector}
    If a real polynomial $p$ satisfies $p(t)\ge1$ for $t\le\eta_0$ and
    $p(h_j)^2\le\Delta$ whenever $h_j\ge\eta_1$, then $p(H)$ is a spectral
    projector for $(H,\eta_0,\eta_1)$ with parameter $\Delta$, with eigenvalues
    $p(h_j)$ on the basis $b_j$.
\end{theorem}

\begin{proof}
    \leanok
    The previous theorem gives $p(H)b_j=p(h_j)b_j$, and the two conditions
    are those of the definition.
\end{proof}

\begin{theorem}[The Chebyshev filter]
    \label{thm:chebyshev_spectral_projector}
    \lean{Submodule.exists_isSpectralProjector_chebyshev}
    \leanok
    \uses{thm:polynomial_spectral_projector, thm:chebyshev_interval_polynomial}
    Let all eigenvalues of $H$ be at most $u$, and let $a<b\le u$. For every
    $n$ there is a real polynomial $p$ of degree $n$ with $p(a)=1$ such that
    $p(H)$ is a spectral projector for $(H,a,b)$ with parameter
    $\bigl(2e^{-2n\sqrt{(b-a)/(u-a)}}\bigr)^2$. This is items 1--3 of
    Theorem~thm:Cheby-AGSP of~\cite[lines 789--853]{Arad2016RigorousRG}, in
    the form required by the error reduction; positivity of $p(H)$ is not
    claimed~\cite{gap:alvv16_chebyshev_filter_positivity}.
\end{theorem}

\begin{proof}
    \leanok
    The normalized interval polynomial satisfies $p(t)\ge1$ for $t\le a$ and
    $\lvert p(t)\rvert\le2e^{-2n\sqrt{(b-a)/(u-a)}}$ on $[b,u]$; apply the
    previous theorem.
\end{proof}

\begin{theorem}[The Chebyshev filter near the normalization point]
    \label{thm:chebyshev_spectral_projector_near}
    \lean{Submodule.exists_isSpectralProjector_chebyshev_one_sub_le}
    \leanok
    \uses{thm:chebyshev_spectral_projector, thm:chebyshev_interval_polynomial_near}
    In the setting of Theorem~\ref{thm:chebyshev_spectral_projector}, the
    polynomial $p$ can be chosen so that, in addition, every eigenvalue
    $h_j\le b$ of $H$ gives an eigenvalue of $p(H)$ with
    \begin{align}
        p(h_j) & \ge1-\frac{n\lvert h_j-a\rvert}{\sqrt{(b-a)(u-a)}}.
        \label{eq:chebyshev_spectral_projector_near}
    \end{align}
    This is item~2 of Theorem~thm:Cheby-AGSP
    of~\cite[lines 822--824]{Arad2016RigorousRG} for eigenvalues near
    $\eta_0=a$, in the corrected form of
    Theorem~\ref{thm:chebyshev_interval_polynomial_near}, on the range
    $h_j\le\eta_1$, which contains the printed range
    $h_j\le\eta_0+\gamma/n$.
\end{theorem}

\begin{proof}
    \leanok
    Use the polynomial of
    Theorem~\ref{thm:chebyshev_interval_polynomial_near}. The previous
    argument gives the spectral projector, and $p(H)b_j=p(h_j)b_j$ with
    (\ref{eq:chebyshev_interval_polynomial_near}) at $x=h_j$ gives
    (\ref{eq:chebyshev_spectral_projector_near}).
\end{proof}
